\section{Conclusion}

\subsection{Summary of Findings}

This study compared the performance of the Louvain and Leiden community detection algorithms
on two biological networks: a bacterial protein–protein interaction network (Bac) 
and a viral protein–protein interaction network (Vir). 
The results, evaluated using the Adjusted Rand Index (ARI), show that both methods identify
highly consistent community structures across both networks. 
This is also reflected in the similar number of communities detected and the comparable modularity scores.

A notable difference observed in this study is that the runtime of
the Leiden algorithm was significantly lower than that of the Louvain algorithm.
This efficiency gain is particularly evident in the larger viral network, 
where Leiden's runtime was approximately an order of magnitude faster than Louvain's.

Overall, the findings suggest that both Louvain and Leiden are 
effective methods for community detection in biological networks.

\subsection{Outlook and Future Work}

This study was originally intended to include an evaluation of the edge-betweenness 
algorithm by Girvan and Newman. However, due to runtime limitations, it was not possible
to include it in the analysis. For future work, it may be worthwhile to explore
ways to improve the computational efficiency of the algorithm. At present, it runs in time 
$O(n^3)$ on sparse graphs, which makes it infeasible for larger networks.

Future research could also further investigate intra-method stability by varying the
resolution parameter in both the Louvain and Leiden algorithms.
Additionally, alternative quality functions beyond modularity could
be explored to address potential limitations such as the resolution limit problem.
Finally, integrating biological validation of detected communities would provide further 
insight into the functional relevance of the identified modular structures.
